\subsection{Recap and the SteaMas scenario}
\label{sec:emas-overview}
\slidesrc{EMAS keynote: 1--15}

The keynote opens with the definitions of Section~\ref{sec:agents-and-mas} (slides 2--5):
the weak definition of agent (situated, autonomous, flexible: reactive, proactive, social),
the strong definition (mentalistic notions: beliefs, desires, intentions, goals, plans), the
four characteristics of a MAS, and the MAS picture of Figure~\ref{fig:mas-structure}
(from Zambonelli and Omicini, 2004). Slide 6 lists further topics of the summer school that
belong to the MAS world: people and organisations; reinforcement learning, deep RL,
multi-agent RL (MARL) and continual learning; coordination, organisations, institutions,
norms and ethics; argumentation; agent-based modelling and simulation.

\begin{nbexample}{the fictional scenario: SteaMas (slides 7--9)}
Dr.\ Joyce, the new mayor of the (steampunk) city of \textbf{SteaMas}, wants to make SteaMas
\emph{green, smart and safe} by ensuring that:
\begin{enumerate}[nosep]
  \item SteaMas citizens and tourists respect well-defined civil/ethical \textbf{norms};
  \item \textbf{no private traffic} exists: citizens and tourists may take advantage of
        different kinds of free public transportation means in the most effective way;
  \item both citizens and tourists may be driven around by a \textbf{personal drone}
        provided by the municipality for free.
\end{enumerate}
Dr.\ Joyce is aware of the existence of agents and MAS, believes that MAS technology can
suit his needs, and contacts us to exploit it. The scenario is resumed in
Section~\ref{sec:steamas}, after the tour of the research directions.
\end{nbexample}

\paragraph{Where EMAS research lives (slides 10--15).} A Google Scholar search for ``agent
oriented software engineering'' returns hundreds of thousands of results. The main venues:
the \textbf{EMAS} workshop (International Workshop on Engineering Multi-Agent Systems),
born from the merger of three predecessors, AOSE (Agent-Oriented Software Engineering),
DALT (Declarative Agent Languages and Technologies) and ProMAS (Programming Multi-Agent
Systems), which started around 2000--2003; and the \emph{Engineering and Analysis of
Multiagent Systems} area of the AAMAS conference (International Conference on Autonomous
Agents and Multiagent Systems), whose topics include requirements and formal specification,
architecture and modelling, verification and validation, programming models and languages,
testing, fault tolerance and reliability, norms and governance, interaction protocols,
BDI-based agents, ethical agents, MAS-based simulations, tools and testbeds, platforms, and
applications from finance to health, agriculture and smart cities.
